\documentclass[11pt]{article}
\usepackage[margin=1in]{geometry}
\usepackage{amsmath,amssymb,booktabs,hyperref,siunitx}
\title{Replication Report: Amplitude Estimation without Phase Estimation\\
\large Suzuki et al., arXiv:1904.10246 (QIP 19, 75, 2020)}
\author{Ollie (subagent), replication run 2026-07-03}
\date{Backfilled 2026-07-05}
\begin{document}
\maketitle

\section{Paper's central claim}
Suzuki et al.\ propose \emph{Maximum Likelihood Amplitude Estimation (MLAE)},
which replaces the QFT-based ancilla register of canonical amplitude estimation
(Brassard--H\o yer--Mosca--Tapp) with a set of independent Grover-power circuits
$A \, Q^{m_k}$ measured directly on the good-state indicator, then combined via
joint MLE over $\{m_k\}$.  Their Fig.\ 2 (right column, $a=1/48$) advertises
three scalings of $\log_{10}\mathrm{RMSE}(a)$ vs.\ $\log_{10} N_q$ (oracle-call
count): classical sampling $-0.50$ (shot noise), LIS $m_k=k$ giving $-0.76$,
and EIS $m_k=2^{k-1}$ giving $-0.95$ (near-Heisenberg).  The headline is that
EIS attains near $1/N$ scaling \emph{without} the ancilla-heavy QPE circuit.

\section{What we exercised}
Independent Qiskit reimplementation.  Single-qubit $A = R_y(2\theta_a)$ and
$Q = -A S_0 A^{\dagger} S_{\chi}$ with $S_{\chi}=Z$, $S_0 = XZX$.  Real
shot-based simulation on \texttt{qiskit-aer.AerSimulator} (not analytic
probabilities).  MLE via fine brute-force grid on $(0,\pi/2)$ plus local
polish, matching the paper's estimator form.  Operating point $a=1/48$,
$N_{\text{shot}}=100$, $100$ trials per $(schedule, M)$ point.  See
\texttt{code/mlae\_replicate.py}, \texttt{report/evidence/results.json}.

\section{Measured slopes (from \texttt{results.json})}
\begin{center}
\begin{tabular}{lrrr}
\toprule
Schedule & Ours & Paper & $\Delta$ \\
\midrule
Classical ($m_k=0$)       & $-0.516$ & $-0.50$ & $+0.016$ \\
LIS ($m_k=k$)             & $-0.727$ & $-0.76$ & $+0.033$ \\
EIS ($m_k=2^{k-1}$)       & $-0.930$ & $-0.95$ & $+0.020$ \\
\bottomrule
\end{tabular}
\end{center}
All three within $\pm 0.04$ of the paper. At matched $N_q \approx 5\times 10^4$
EIS delivers RMSE $\approx 5\times 10^{-5}$ vs.\ classical extrapolation
$\approx 6\times 10^{-4}$ (\textbf{$\sim 10\times$ tighter}), exactly the
``quantum advantage without QFT'' the paper advertises.

\section{Verdict}
\textbf{REPLICATED.}  Headline exercised: (i) independent from-scratch Qiskit
implementation of $A$, $Q$, and the joint MLE; (ii) all three schedules'
scaling slopes match the paper's Fig.\ 2 within $\pm 0.04$; (iii) EIS's
approach to Heisenberg scaling ($-0.93$ vs $-0.95$ theoretical target,
clearly steeper than the classical $-0.50$) directly reproduced; (iv) the
$\sim 10\times$ error reduction at matched query budget confirmed
quantitatively.

\section{Honest critique --- what we did NOT do}
\begin{itemize}
  \item \textbf{No head-to-head against canonical QAE.} We did not implement
        Brassard--H\o yer--Mosca--Tapp (QPE-based) to numerically confirm the
        \emph{gate-depth / ancilla-count} savings claim; only the estimator
        scaling was checked.  Depth-vs-accuracy tradeoff is asserted, not
        measured.
  \item \textbf{Single amplitude only.} We reproduced the $a=1/48$ operating
        point.  Behavior across the range where MLE can pick a wrong branch
        (near $\theta_a$ close to $\pi/4$, or very small $a$ with limited
        shot budget) was not swept; identifiability failure modes untested.
  \item \textbf{Noiseless simulator.} \texttt{qiskit-aer} with zero noise.
        The paper's practical implication for near-term hardware --- that
        MLAE's shallower circuits are more noise-robust than QPE --- was
        not tested.  This is arguably the most important open question and
        is a limitation we flag prominently in \texttt{failure\_analysis.md}
        and the open-questions list.
  \item \textbf{MLE convergence not stress-tested.} Grid + polish worked at
        our shot counts, but the paper does not, and neither did we,
        exhaustively characterize local-minimum failure rates for small
        $N_{\text{shot}}$ or adversarial schedules.
  \item \textbf{No comparison to modern successors.} Iterative QAE (Grinko
        et al.), Faster QAE (Nakaji), and shot-optimized variants have since
        superseded MLAE on several axes.  Not benchmarked here.
  \item \textbf{Code fix (transpile cache) risk.} We added memoization of
        \texttt{transpile()} keyed on $(\theta_a, m)$.  This changes no
        physics --- the transpiled circuit is a pure function of
        $(\theta_a, m)$ --- but it is a modification of the pre-existing
        draft; slope numbers should be identical, which they were.
\end{itemize}

\section{Files}
\begin{itemize}
  \item \texttt{code/mlae\_replicate.py} --- circuit build, sampling, MLE, slope fit
  \item \texttt{report/evidence/results.json} --- raw per-point RMSE + slopes
  \item \texttt{logs/main.log} --- run stdout
  \item \texttt{report/REPORT.md} --- prose version of this report
  \item \texttt{report/open\_questions.json} --- five open questions with next steps
  \item \texttt{report/failure\_analysis.md} --- limitations and negative results
  \item \texttt{report/workflow.md} --- replication workflow
  \item \texttt{report/artifacts\_summary.md} --- artifact inventory
\end{itemize}
\end{document}
